% ==========================================================
% Section: Introduction
% ==========================================================

\section{Introduction}
\label{sec:introduction}

\dropcapstart{W}ith the rapid advancement of modern Internet of Things (IoT) ecosystems and cyber-physical infrastructure, automated Radio Frequency Identification (RFID) systems have gained immense traction across supply chains, smart manufacturing, and logistics~\cite{Industry5, RFID-application1}. Fast and reliable identification of dense dynamic tags is essential for maintaining high throughput and mission-critical tracking fidelity.

\subsection{Research Background and Motivation}
In industrial environments, tags move dynamically across interrogation zones under stochastic residence intervals~\cite{RFID-in-IOT}. Conventional anti-collision protocols based on static frame-slotted ALOHA or deterministic tree splitting encounter severe communication latency and tag starvation under heavy workload surges~\cite{EPC-C1-Gen2, TAD}. Consequently, architecting a unified, modular framework that decouples the scientific algorithmic representation from venue-specific typography is of substantial benefit to reproducible research and rapid dissemination.

\subsection{Primary Contributions}
The major contributions of this work are three-fold:
\begin{itemize}
    \item We design a decoupled multi-publisher template architecture that isolates manuscript content from target venue typographical dependencies.
    \item We seamlessly integrate official journal standards (including IEEE Transactions, Elsevier CAS, and ACM SIGCONF) with full support for authentic bibliographic databases.
    \item We empirically validate seamless template switching without modifying the underlying technical discourse or experimental data.
\end{itemize}
